\begin{filecontents}{references.bib}

@article{Awang2025,
  author = {Awang, Y. and Yusop, F. and Danaee, M.},
  title = {Current practices and future direction of AI in mathematics education},
  journal = {Education and Information Technologies},
  year = {2025}
}

@article{Walkington2025,
  author = {Walkington, Candace},
  title = {The implications of generative AI for mathematics education},
  journal = {Journal of Mathematics Education},
  year = {2025}
}

@article{Jacobs2022,
  author = {Jacobs, J.},
  title = {Artificial Intelligence and Mathematics Teaching},
  journal = {Educational Studies in Mathematics},
  year = {2022}
}

@article{Rezat2021,
  author = {Rezat, S.},
  title = {Digital Technologies in Mathematics Education},
  journal = {International Journal of STEM Education},
  year = {2021}
}

@article{VillegasCh2025,
  author = {Villegas-Ch, W.},
  title = {Intelligent Tutoring Systems in Mathematics Education},
  journal = {Education Sciences},
  year = {2025}
}

@article{Feng2025,
  author = {Feng, Y.},
  title = {AI Applications in School Mathematics},
  journal = {Computers and Education},
  year = {2025}
}

@article{Demszky2025,
  author = {Demszky, D.},
  title = {Generative Artificial Intelligence in Mathematics Education},
  journal = {Computers and Education},
  year = {2025}
}

@article{Steinbach2025,
  author = {Steinbach, M.},
  title = {Risks of Generative AI in Education},
  journal = {AI and Society},
  year = {2025}
}

@article{Kock2025192,
  author  = {Kock, Zeger-Jan and Salinas-Hernández, Ulises and Pepin, Birgit},
  title   = {Engineering Students’ Initial Use Schemes of ChatGPT as an Instrument for Learning},
  journal = {Digital Experiences in Mathematics Education},
  year    = {2025},
  volume   = {11},
  number   = {1},
  pages    = {192--218},
  doi      = {10.1007/s40751-025-00169-w}
}

@ARTICLE{Taani20251748,
 author = {Taani, Osama and Alabidi, Suzan},
 title = {ChatGPT in education: benefits and challenges of ChatGPT for mathematics and science teaching practices},
 year = {2025},
 journal = {International Journal of Mathematical Education in Science and Technology},
 volume = {56},
 number = {9},
 pages = {1748--1777},
 doi = {10.1080/0020739X.2024.2357341}
}

@article{UNESCO2023,
  author  = {{UNESCO}},
  title   = {Guidance for Generative AI in Education and Research},
  journal = {UNESCO Policy Document},
  year    = {2023},
  publisher = {UNESCO},
  address = {Paris}
}

@article{AIAct2024,
  author  = {{European Parliament}},
  title   = {Artificial Intelligence Act},
  journal = {Official Journal of the European Union},
  year    = {2024},
  publisher = {European Union}
}

@article{ConceptAI2020,
  author  = {{Кабінет Міністрів України}},
  title   = {Концепція розвитку штучного інтелекту в Україні},
  journal = {Урядовий нормативно-правовий документ},
  year    = {2020},
  note    = {Розпорядження Кабінету Міністрів України №1556-р}
}

\end{filecontents}

\documentclass[14pt,a4paper]{extarticle}

\usepackage{polyglossia}
\setdefaultlanguage{ukrainian}

\usepackage{fontspec}
\setmainfont{DejaVu Serif}
\newfontfamily\cyrillicfont{DejaVu Serif}
\setmonofont{DejaVu Sans Mono}
\newfontfamily\cyrillicfonttt{DejaVu Sans Mono}

\usepackage[a4paper,
left=30mm,
right=15mm,
top=20mm,
bottom=20mm]{geometry}

\usepackage{setspace}
\onehalfspacing

\usepackage{indentfirst}

\usepackage{float}
\usepackage{longtable}

\setlength{\parindent}{1.25cm}

\usepackage{hyperref}
\usepackage{csquotes}

\usepackage[
backend=biber,
style=authoryear, %apa,
]{biblatex}

\addbibresource{references.bib}

\usepackage{microtype}
\emergencystretch=3em

\begin{document}

% =========================
% ТИТУЛЬНА СТОРІНКА
% =========================

\begin{titlepage}

\begin{center}

МІНІСТЕРСТВО ОСВІТИ І НАУКИ УКРАЇНИ

\vspace{1cm}

КРИВОРІЗЬКИЙ ДЕРЖАВНИЙ ПЕДАГОГІЧНИЙ УНІВЕРСИТЕТ

\vspace{0.5cm}

Фізико-математичний факультет

\vspace{0.5cm}

Кафедра математики та методики її навчання

\vspace{5cm}

{\Large\textbf{КУРСОВА РОБОТА}}

\vspace{0.5cm}

на тему:

\vspace{0.5cm}

{\Large\textbf{
ЗАСТОСУВАННЯ ШТУЧНОГО ІНТЕЛЕКТУ
В РОБОТІ ВЧИТЕЛЯ МАТЕМАТИКИ
}}

\end{center}

\vspace{3cm}

\begin{flushright}

Виконала: \\
студентка групи МІ-22 \\
Осипчук Валерія Олександрівна

\vspace{1cm}

Науковий керівник: \\
Семеріков Сергій Олексійович

\end{flushright}

\vfill

\begin{center}
Кривий Ріг — 2026
\end{center}

\end{titlepage}

% =========================
% ЗМІСТ
% =========================

\tableofcontents

\newpage

% =========================
% ВСТУП
% =========================

\section*{ВСТУП}
\addcontentsline{toc}{section}{ВСТУП}

Актуальність теми дослідження. Стрімкий розвиток цифрових технологій та активне впровадження систем штучного інтелекту в різні сфери суспільного життя зумовлюють необхідність модернізації освітнього процесу та пошуку нових підходів до організації навчання. Одним із пріоритетних напрямів цифрової трансформації освіти є використання технологій штучного інтелекту (Artificial Intelligence, AI), які відкривають широкі можливості для персоналізації навчання, автоматизації рутинних педагогічних завдань, аналізу навчальних результатів та створення адаптивного освітнього середовища \cite{Awang2025,Walkington2025}.

Особливої актуальності зазначена проблема набуває в математичній освіті. Математика є однією з фундаментальних навчальних дисциплін, успішне засвоєння якої потребує систематичного контролю навчальних досягнень, індивідуального підходу до учнів та своєчасного надання зворотного зв'язку. Сучасні AI-технології здатні забезпечувати підтримку цих процесів шляхом автоматизованого оцінювання, створення математичних завдань, аналізу типових помилок учнів, формування індивідуальних освітніх траєкторій та організації адаптивного навчання \cite{VillegasCh2025,Feng2025}.

Упродовж останніх років особливого поширення набули генеративні моделі штучного інтелекту, зокрема ChatGPT, Claude, Gemini, Microsoft Copilot та інші великі мовні моделі (Large Language Models). Їх використання у професійній діяльності педагога дозволяє значно скоротити час на підготовку навчальних матеріалів, створення тестових завдань, розроблення диференційованих вправ та організацію індивідуальної роботи учнів \cite{Demszky2025}. Водночас впровадження генеративного AI супроводжується низкою викликів, пов'язаних із достовірністю отриманих результатів, дотриманням принципів академічної доброчесності, захистом персональних даних та необхідністю розвитку критичного мислення учнів \cite{Steinbach2025}.

Проблеми цифровізації освіти та використання інформаційно-комунікаційних технологій у навчальному процесі досліджували В. Ю. Биков, Н. В. Морзе, О. М. Спірін, С. О. Семеріков та інші вітчизняні науковці. Теоретичні та практичні аспекти використання штучного інтелекту в освіті висвітлено у працях Y. Awang, C. Walkington, W. Villegas-Ch, D. Demszky, Y. Feng, B. Pepin, O. Taani та інших зарубіжних дослідників. Проведений аналіз наукових джерел свідчить про значний потенціал технологій штучного інтелекту для вдосконалення математичної освіти, проте питання їх ефективного та педагогічно доцільного використання в роботі вчителя математики потребують подальшого вивчення.

Актуальність зазначеної проблеми, її теоретичне та практичне значення для сучасної математичної освіти зумовили вибір теми курсової роботи: \textbf{«Застосування штучного інтелекту в роботі вчителя математики»}.

\textbf{Об'єкт дослідження} – процес навчання математики в закладах загальної середньої освіти.

\textbf{Предмет дослідження} – використання технологій штучного інтелекту в професійній діяльності вчителя математики.

\textbf{Мета дослідження} – теоретично обґрунтувати особливості використання технологій штучного інтелекту в роботі вчителя математики та визначити перспективні напрями їх застосування в освітньому процесі.

Для досягнення поставленої мети необхідно виконати такі завдання:

\begin{enumerate}
\item проаналізувати сучасні наукові підходи до використання штучного інтелекту в освіті;
\item охарактеризувати основні напрями застосування AI-технологій у математичній освіті;
\item дослідити можливості генеративного штучного інтелекту та інтелектуальних навчальних систем у професійній діяльності вчителя математики;
\item визначити переваги, ризики та обмеження використання AI в освітньому процесі;
\item узагальнити сучасний досвід використання AI на уроках математики;
\item сформулювати методичні рекомендації щодо ефективного використання технологій штучного інтелекту в роботі вчителя математики.
\end{enumerate}

\textbf{Методи дослідження}: аналіз і синтез наукових джерел, узагальнення та систематизація результатів сучасних досліджень, порівняльний аналіз підходів до використання штучного інтелекту в освіті, аналіз практичного досвіду застосування AI-технологій у математичній освіті.

\textbf{Практичне значення дослідження} полягає в узагальненні сучасних підходів до використання технологій штучного інтелекту в роботі вчителя математики та розробленні рекомендацій щодо їх застосування під час підготовки навчальних матеріалів, створення математичних завдань, організації диференційованого навчання та оцінювання навчальних досягнень учнів.

Теоретичну основу дослідження становлять праці вітчизняних і зарубіжних учених, присвячені цифровізації освіти, використанню штучного інтелекту в навчальному процесі та особливостям математичної освіти в умовах розвитку сучасних цифрових технологій. Інформаційну базу дослідження складають наукові статті, проіндексовані в міжнародній наукометричній базі Scopus, а також сучасні публікації з проблематики використання штучного інтелекту в освіті.

Структура курсової роботи зумовлена метою та завданнями дослідження. Робота складається зі вступу, трьох розділів, висновків, списку використаних джерел та додатків.

% =========================
% РОЗДІЛ 1
% =========================

\section{ТЕОРЕТИЧНІ ОСНОВИ ВИКОРИСТАННЯ ШТУЧНОГО ІНТЕЛЕКТУ В МАТЕМАТИЧНІЙ ОСВІТІ}

\subsection{Штучний інтелект у сучасній освіті}

Штучний інтелект є одним із найдинамічніших напрямів розвитку сучасних цифрових технологій. Упродовж останніх років AI активно інтегрується в освітній процес, створюючи нові можливості для персоналізації навчання, автоматизації оцінювання та вдосконалення педагогічної діяльності.

Дослідники зазначають, що сучасні AI-технології сприяють переходу від традиційної моделі навчання до адаптивного освітнього середовища, орієнтованого на індивідуальні потреби учнів.

У сучасній освіті активно використовуються такі напрями AI:

\begin{itemize}
\item генеративний штучний інтелект;
\item інтелектуальні навчальні системи;
\item адаптивні освітні платформи;
\item автоматизоване оцінювання;
\item AI-чатботи;
\item системи аналізу навчальних даних.
\end{itemize}

Використання AI дозволяє:

\begin{itemize}
\item підвищити ефективність навчального процесу;
\item забезпечити індивідуальний підхід;
\item скоротити час на виконання рутинних педагогічних завдань;
\item покращити доступність освітніх ресурсів.
\end{itemize}

Важливим напрямом розвитку AI в освіті є використання генеративного штучного інтелекту. Генеративні моделі здатні створювати навчальні матеріали, адаптувати тексти до рівня підготовки учнів та формувати персоналізований освітній контент.

Поява ChatGPT, Claude, Gemini та інших LLM-систем суттєво вплинула на освітню сферу. Такі технології дозволяють:

\begin{itemize}
\item генерувати пояснення;
\item створювати вправи;
\item формувати тести;
\item аналізувати відповіді учнів;
\item організовувати інтерактивну підтримку навчання.
\end{itemize}

Науковці наголошують, що AI не повинен розглядатися як повноцінна заміна педагога. Основна роль учителя полягає у педагогічному супроводі, розвитку критичного мислення та формуванні навчальної мотивації учнів.

\subsection{Особливості використання AI у навчанні математики}

Математична освіта є однією з галузей, де AI-технології демонструють високий потенціал ефективності. Це пояснюється можливістю автоматизованого аналізу відповідей, адаптивного формування завдань та персоналізації освітнього процесу.

Одним із ключових напрямів є використання інтелектуальних навчальних систем (ITS). Такі системи аналізують навчальні результати учнів та адаптують зміст завдань відповідно до рівня підготовки.

Інтелектуальні системи забезпечують:

\begin{itemize}
\item миттєвий зворотний зв’язок;
\item виявлення типових помилок;
\item адаптацію складності завдань;
\item підтримку індивідуальної траєкторії навчання.
\end{itemize}

Значного поширення набуває використання генеративного AI у процесі навчання математики. ChatGPT та інші LLM-моделі використовуються для:

\begin{itemize}
\item пояснення математичних понять;
\item створення прикладів;
\item генерації математичних задач;
\item підготовки тестових завдань;
\item створення диференційованих вправ.
\end{itemize}

Дослідження показують, що використання ChatGPT може позитивно впливати на навчальну мотивацію та активність учнів. Водночас науковці підкреслюють необхідність перевірки AI-відповідей учителем.

Важливим напрямом є автоматизований аналіз помилок учнів. Системи на основі AI дозволяють швидко визначати типові помилки та формувати персоналізовані рекомендації щодо їх виправлення.

Перспективним напрямом є використання AI для розв’язування математичних текстових задач. Генеративні моделі можуть пояснювати алгоритм розв’язання та демонструвати покрокове виконання задач.

\subsection{Роль учителя математики в умовах розвитку AI}

Розвиток AI-технологій змінює професійну діяльність учителя математики. Педагог поступово переходить від ролі основного джерела знань до ролі фасилітатора навчального процесу.

Сучасний учитель математики повинен:

\begin{itemize}
\item володіти цифровими компетентностями;
\item уміти використовувати AI-сервіси;
\item оцінювати достовірність AI-відповідей;
\item формувати критичне мислення учнів;
\item забезпечувати академічну доброчесність.
\end{itemize}

Модель AI-TPACK розглядається як одна з перспективних концепцій підготовки сучасного педагога. Вона поєднує:

\begin{itemize}
\item предметні знання;
\item педагогічні знання;
\item технологічні компетентності;
\item навички використання AI.
\end{itemize}

AI може виступати ефективним інструментом підтримки професійної діяльності вчителя:

\begin{itemize}
\item автоматизація підготовки матеріалів;
\item створення тестів;
\item генерація задач;
\item підготовка презентацій;
\item аналіз результатів навчання.
\end{itemize}

Разом із тим ключова роль педагога зберігається, оскільки саме вчитель забезпечує:

\begin{itemize}
\item педагогічну взаємодію;
\item емоційну підтримку;
\item розвиток критичного мислення;
\item виховний компонент навчання.
\end{itemize}

\subsection{Бар'єри та ризики використання AI в освіті}

Попри значні переваги, використання AI в освіті супроводжується низкою ризиків та обмежень.

Однією з основних проблем є академічна недоброчесність. Генеративний AI може використовуватися учнями для автоматичного виконання домашніх завдань без самостійного аналізу матеріалу.

Іншим ризиком є можливість генерації недостовірної інформації. LLM-моделі іноді формують помилкові математичні пояснення або некоректні відповіді.

До основних ризиків використання AI належать:

\begin{itemize}
\item надмірна залежність від технологій;
\item зниження рівня самостійного мислення;
\item ризики конфіденційності даних;
\item цифрова нерівність;
\item недостатня AI-компетентність педагогів.
\end{itemize}

Проблема цифрової нерівності є особливо актуальною для закладів освіти з недостатнім технічним забезпеченням.

Науковці також наголошують на необхідності формування етичних принципів використання AI в освіті. Учителі повинні:

\begin{itemize}
\item контролювати використання AI;
\item перевіряти достовірність відповідей;
\item навчати учнів критично оцінювати інформацію;
\item дотримуватися принципів академічної доброчесності.
\end{itemize}

Отже, використання штучного інтелекту в математичній освіті має значний потенціал, однак потребує педагогічно обґрунтованого та відповідального впровадження.

\subsection*{Висновки до розділу 1}

У першому розділі було проаналізовано теоретичні основи використання штучного інтелекту в математичній освіті.

Встановлено, що AI-технології активно інтегруються в сучасний освітній процес та створюють нові можливості для персоналізації навчання, автоматизації оцінювання та підтримки професійної діяльності вчителя математики.

Визначено, що AI може ефективно використовуватися для створення математичних завдань, аналізу помилок учнів та організації індивідуальної освітньої траєкторії.

Разом із тим ефективність використання AI залежить від рівня цифрової компетентності педагога, дотримання принципів академічної доброчесності та педагогічно обґрунтованого використання цифрових технологій.0


% =========================
% РОЗДІЛ 2
% =========================

\section{ПРАКТИЧНІ АСПЕКТИ ВИКОРИСТАННЯ AI}

Генеративний AI може використовуватися для створення математичних задач, тестів та диференційованих вправ \parencite{Walkington2025}.

ChatGPT та інші генеративні моделі можуть допомагати вчителю під час підготовки навчальних матеріалів, пояснення складних тем та організації індивідуальної роботи учнів.


\subsection{Використання генеративного AI у підготовці до уроків}

Сучасні генеративні AI-системи відкривають нові можливості для організації професійної діяльності вчителя математики. Найбільш поширеними інструментами є ChatGPT, Claude, Gemini, Copilot та інші LLM-моделі, здатні генерувати текстовий контент, математичні пояснення, завдання та навчальні матеріали \cite{Walkington2025,Kock2025192}.

Одним із головних напрямів використання генеративного AI є підготовка до уроків. Учитель математики може використовувати AI для:
\begin{itemize}
\item створення конспектів уроків;
\item генерації математичних задач;
\item підготовки тестових завдань;
\item створення диференційованих вправ;
\item формування прикладів розв’язування задач;
\item створення презентацій;
\item підготовки домашніх завдань.
\end{itemize}

Науковці зазначають, що використання AI дозволяє суттєво скоротити час підготовки до занять \\cite{Awang2025}. Завдяки генеративним моделям учитель може швидко адаптувати навчальний матеріал відповідно до рівня підготовки учнів.

Особливо корисною є можливість створення диференційованих математичних вправ. Генеративний AI може формувати:
\begin{itemize}
\item базові вправи;
\item задачі підвищеної складності;
\item практикоорієнтовані задачі;
\item текстові задачі;
\item творчі математичні завдання.
\end{itemize}

Дослідження \cite{Awang2025} показують, що AI допомагає педагогам створювати більш структуровані та варіативні навчальні матеріали.

Водночас учитель повинен критично оцінювати результати роботи AI. Генеративні моделі можуть:
\begin{itemize}
\item допускати математичні помилки;
\item створювати некоректні пояснення;
\item формувати задачі з неоднозначними умовами.
\end{itemize}

Тому використання AI повинно здійснюватися під постійним контролем педагога.

\subsection{Використання ChatGPT для створення математичних задач}

ChatGPT є одним із найбільш популярних AI-інструментів у сфері освіти. Його використання у математичній освіті дозволяє:
\begin{itemize}
\item генерувати математичні вправи;
\item створювати приклади розв’язування;
\item пояснювати математичні поняття;
\item адаптувати складність матеріалу;
\item створювати індивідуальні завдання.
\end{itemize}

Важливою перевагою ChatGPT є можливість швидкого створення великої кількості варіантів математичних задач. Наприклад, учитель може сформулювати запит щодо генерації:
\begin{itemize}
\item задач на дроби;
\item рівнянь;
\item геометричних задач;
\item текстових задач;
\item логічних вправ.
\end{itemize}

Генеративний AI також може створювати задачі різного рівня складності для реалізації диференційованого навчання \cite{Walkington2025}.

Приклад використання ChatGPT:
\begin{quote}
«Створи 5 текстових задач на відсотки для учнів 6 класу з практичним змістом».
\end{quote}

У результаті AI генерує набір задач, які вчитель може адаптувати відповідно до навчальної мети уроку.

Науковці \cite{Taani20251748} підкреслюють, що використання ChatGPT позитивно впливає на мотивацію учнів завдяки інтерактивності та можливості швидкого отримання пояснень.

Разом із тим існують певні ризики:
\begin{itemize}
\item учні можуть бездумно копіювати відповіді;
\item можливе зниження рівня самостійного мислення;
\item AI може генерувати помилкові розв’язання.
\end{itemize}

Саме тому педагог повинен формувати в учнів навички критичного аналізу AI-відповідей.

\subsection{AI-системи для автоматизованого оцінювання}

Одним із перспективних напрямів є використання AI для автоматизованого оцінювання навчальних результатів учнів \cite{Walkington2025,Jacobs2022}.

AI-системи дозволяють:
\begin{itemize}
\item автоматично перевіряти тести;
\item аналізувати типові помилки;
\item формувати персоналізований зворотний зв’язок;
\item оцінювати динаміку навчальних досягнень.
\end{itemize}

Особливу ефективність такі системи демонструють у процесі:
\begin{itemize}
\item перевірки математичних вправ;
\item оцінювання тестових завдань;
\item аналізу помилок у розв’язаннях;
\item моніторингу навчального прогресу.
\end{itemize}

Автоматизований зворотний зв’язок дозволяє учням швидше виявляти помилки та коригувати власну діяльність \cite{Rezat2021}.

Перевагами AI-оцінювання є:
\begin{itemize}
\item економія часу вчителя;
\item об’єктивність оцінювання;
\item швидке отримання результатів;
\item персоналізація рекомендацій.
\end{itemize}

Проте повна автоматизація оцінювання має певні обмеження. AI не завжди здатний:
\begin{itemize}
\item оцінити нестандартне мислення;
\item врахувати творчий підхід;
\item визначити логіку міркувань учня.
\end{itemize}

Тому остаточне оцінювання повинно залишатися за педагогом.

\subsection{Інтелектуальні навчальні системи в математичній освіті}

Інтелектуальні навчальні системи (ITS) є одним із найперспективніших напрямів використання AI у математичній освіті \cite{VillegasCh2025}.

Такі системи забезпечують:
\begin{itemize}
\item адаптивне навчання;
\item аналіз навчальних результатів;
\item індивідуалізацію освітнього процесу;
\item підтримку самостійного навчання.
\end{itemize}

ITS-системи використовують алгоритми AI для аналізу:
\begin{itemize}
\item успішності учнів;
\item типових помилок;
\item темпу навчання;
\item рівня складності завдань.
\end{itemize}

На основі отриманих даних система автоматично адаптує навчальний матеріал.

Дослідження \cite{Feng2025} демонструють, що ITS позитивно впливають на:
\begin{itemize}
\item навчальну мотивацію;
\item рівень математичних досягнень;
\item самостійність учнів;
\item ефективність навчання.
\end{itemize}

Особливу ефективність ITS демонструють під час:
\begin{itemize}
\item вивчення алгебри;
\item формування навичок розв’язування задач;
\item розвитку математичного мислення;
\item підготовки до тестування.
\end{itemize}

Разом із тим впровадження ITS потребує:
\begin{itemize}
\item належного технічного забезпечення;
\item цифрової компетентності педагогів;
\item методичної підготовки вчителів;
\item контролю достовірності AI-рекомендацій.
\end{itemize}

\subsection{Порівняльний аналіз AI-інструментів у роботі вчителя математики}

Для ефективного використання AI у професійній діяльності важливо враховувати функціональні можливості різних\\
AI-інструментів.

\begin{table}[H]
\centering
\caption{Порівняльна характеристика AI-інструментів у математичній освіті}
\label{tab:ai_tools}

\renewcommand{\arraystretch}{1.3}

\begin{tabular}{|p{2.8cm}|p{3.8cm}|p{3.8cm}|p{3.8cm}|}
\hline

\textbf{AI-інструмент} &
\textbf{Основні можливості} &
\textbf{Переваги} &
\textbf{Ризики} \\

\hline

ChatGPT &
генерація текстів, задач, пояснень &
швидкість роботи, доступність, інтерактивність &
можливі помилки у відповідях, академічна недоброчесність \\

\hline

Claude &
структурований аналіз текстів, створення матеріалів &
якісне пояснення та узагальнення інформації &
потребує перевірки результатів \\

\hline

Gemini &
мультимодальна підтримка навчання &
інтеграція з сервісами Google &
можливі неточності у відповідях \\

\hline

ITS-системи &
адаптивне навчання та оцінювання &
персоналізація навчання &
потреба у технічному забезпеченні \\

\hline

Photomath &
покрокове розв’язування задач &
візуалізація розв’язання &
ризик механічного копіювання \\

\hline

Wolfram Alpha &
математичні обчислення та побудова графіків &
висока точність обчислень &
обмеження у поясненні логіки \\

\hline

\end{tabular}
\end{table}

Отже, AI-технології мають значний потенціал для вдосконалення професійної діяльності вчителя математики, однак їх використання повинно бути педагогічно доцільним та контрольованим.

\subsection*{Висновки до розділу 2}

У другому розділі було досліджено практичні аспекти використання AI в роботі вчителя математики. Встановлено, що генеративний штучний інтелект та інтелектуальні навчальні системи можуть ефективно використовуватися для створення навчальних матеріалів, автоматизованого оцінювання та організації адаптивного навчання.

Проведений аналіз показав, що використання AI сприяє персоналізації освітнього процесу, економії часу педагога та підвищенню навчальної мотивації учнів.

Разом із тим визначено, що ефективність використання AI залежить від цифрової компетентності педагога, педагогічного контролю та дотримання принципів академічної доброчесності.
```


% =========================
% РОЗДІЛ 3
% =========================

\section{МЕТОДИЧНІ ОСОБЛИВОСТІ ВИКОРИСТАННЯ AI}

Ефективність використання AI залежить від цифрової компетентності вчителя та педагогічного контролю.

Використання AI не повинно замінювати професійну діяльність педагога, а має виступати інструментом підтримки освітнього процесу.



% =========================
% РОЗДІЛ 3
% =========================


\subsection{Використання AI-технологій на різних етапах уроку математики}

Штучний інтелект може ефективно використовуватися на різних етапах уроку математики. AI-технології дозволяють зробити освітній процес більш інтерактивним, адаптивним та індивідуалізованим.

На етапі мотивації навчальної діяльності вчитель може використовувати генеративний AI для:
\begin{itemize}
\item створення проблемних ситуацій;
\item формування практикоорієнтованих задач;
\item підготовки інтерактивних математичних вікторин;
\item створення візуалізацій математичних понять.
\end{itemize}

Наприклад, ChatGPT може генерувати життєві математичні ситуації, що сприяють підвищенню зацікавленості учнів у вивченні теми.

Під час пояснення нового матеріалу AI може використовуватися для:
\begin{itemize}
\item створення покрокових пояснень;
\item побудови графіків;
\item візуалізації математичних моделей;
\item формування прикладів різного рівня складності.
\end{itemize}

Особливо ефективним є використання AI у процесі пояснення складних математичних понять. Генеративні моделі здатні адаптувати пояснення відповідно до віку та рівня підготовки учнів \cite{Walkington2025}.

На етапі закріплення знань AI-технології дозволяють:
\begin{itemize}
\item автоматично генерувати вправи;
\item створювати варіативні завдання;
\item організовувати адаптивне тренування;
\item забезпечувати миттєвий зворотний зв’язок.
\end{itemize}

Важливою перевагою AI є можливість організації диференційованого навчання. Учитель може створювати індивідуальні набори вправ відповідно до рівня навчальних досягнень учнів.

На етапі контролю знань AI може використовуватися для:
\begin{itemize}
\item автоматизованого тестування;
\item аналізу помилок;
\item оцінювання результатів;
\item моніторингу навчального прогресу.
\end{itemize}

Дослідження \cite{Demszky2025} підтверджують, що AI-системи можуть суттєво підвищувати ефективність зворотного зв’язку та підтримувати педагогічну діяльність учителя.

\subsection{Методика використання генеративного AI на уроках математики}

Використання генеративного AI вимагає методично обґрунтованого підходу. Ефективність застосування AI залежить від:
\begin{itemize}
\item навчальної мети;
\item рівня підготовки учнів;
\item типу математичного матеріалу;
\item педагогічної майстерності вчителя.
\end{itemize}

Під час використання ChatGPT або інших генеративних моделей учитель повинен:
\begin{itemize}
\item формулювати чіткі запити;
\item перевіряти правильність AI-відповідей;
\item адаптувати матеріал до вікових особливостей учнів;
\item забезпечувати розвиток критичного мислення.
\end{itemize}

Приклад ефективного AI-запиту:
\begin{quote}
«Створи 3 задачі на відсотки для учнів 7 класу з практичним змістом та покроковим поясненням».
\end{quote}

Генеративний AI може використовуватися для:
\begin{itemize}
\item створення дидактичних матеріалів;
\item підготовки тестів;
\item формування домашніх завдань;
\item підготовки математичних досліджень;
\item створення інтерактивних вправ.
\end{itemize}

Водночас учитель повинен враховувати можливі недоліки AI:
\begin{itemize}
\item математичні помилки;
\item формальні пояснення;
\item відсутність педагогічного контексту;
\item можливість некоректних відповідей.
\end{itemize}

Тому генеративний AI слід розглядати як допоміжний інструмент, а не як повноцінну заміну педагогічної діяльності.

\subsection{Розвиток критичного мислення учнів під час використання AI}

Одним із ключових завдань сучасної математичної освіти є формування критичного мислення учнів. Використання AI-технологій створює як нові можливості, так і нові виклики у цьому напрямі.

Учитель математики повинен навчати учнів:
\begin{itemize}
\item перевіряти достовірність AI-відповідей;
\item аналізувати математичні пояснення;
\item знаходити помилки в AI-розв’язаннях;
\item самостійно обґрунтовувати результати.
\end{itemize}

Ефективним методичним прийомом є аналіз помилкових AI-відповідей. Учитель може спеціально демонструвати учням некоректні розв’язання, згенеровані AI, та організовувати їх колективне обговорення.

Такий підхід сприяє:
\begin{itemize}
\item розвитку логічного мислення;
\item формуванню математичної грамотності;
\item розвитку навичок аналізу інформації;
\item підвищенню цифрової компетентності.
\end{itemize}

Науковці \cite{Steinbach2025} підкреслюють, що учні повинні усвідомлювати обмеження генеративного AI та вміти критично оцінювати отримані результати.

\subsection{Методичні рекомендації щодо використання AI в роботі вчителя математики}

На основі проведеного аналізу можна сформулювати основні методичні рекомендації щодо використання AI в професійній діяльності вчителя математики.

\textbf{1. Використовувати AI як допоміжний інструмент.}

AI повинен підтримувати педагогічну діяльність учителя, а не замінювати її. Остаточний контроль навчального процесу має залишатися за педагогом.

\textbf{2. Перевіряти правильність AI-відповідей.}

Учитель повинен аналізувати:
\begin{itemize}
\item математичну коректність;
\item логічність пояснень;
\item відповідність навчальній програмі;
\item доступність матеріалу для учнів.
\end{itemize}

\textbf{3. Формувати критичне мислення учнів.}

Учні повинні:
\begin{itemize}
\item аналізувати AI-відповіді;
\item перевіряти результати;
\item аргументувати власну думку;
\item уникати механічного копіювання.
\end{itemize}

\textbf{4. Дотримуватися принципів академічної доброчесності.}

Учитель має пояснювати учням правила відповідального використання AI та запобігати порушенням академічної доброчесності.

\textbf{5. Постійно підвищувати цифрову компетентність.}

Сучасний педагог повинен:
\begin{itemize}
\item володіти навичками роботи з AI;
\item знати можливості сучасних AI-сервісів;
\item орієнтуватися в етичних питаннях використання AI.
\end{itemize}

\subsection{Порівняльний аналіз переваг та ризиків використання AI на уроках математики}

\begin{table}[H]
\centering

\caption{Переваги та ризики використання AI в математичній освіті}
\label{tab:advantages_risks}

\renewcommand{\arraystretch}{1.4}

\begin{tabular}{|p{6cm}|p{6cm}|}
\hline

\textbf{Переваги} &
\textbf{Ризики} \\

\hline

Персоналізація навчання &
залежність учнів від AI \\

\hline

Швидке створення навчальних матеріалів &
можливі математичні помилки \\

\hline

Автоматизований зворотний зв’язок &
порушення академічної доброчесності \\

\hline

Підвищення мотивації учнів &
зниження рівня самостійного мислення \\

\hline

Економія часу вчителя &
недостатня цифрова компетентність педагогів \\

\hline

Індивідуалізація освітнього процесу &
ризики конфіденційності даних \\

\hline

Адаптивне навчання &
цифрова нерівність \\

\hline

\end{tabular}

\end{table}

Таким чином, ефективність використання AI у математичній освіті залежить від педагогічно обґрунтованого підходу, цифрової компетентності вчителя та відповідального використання AI-технологій.

\subsection*{Висновки до розділу 3}

У третьому розділі було досліджено методичні особливості використання AI на уроках математики. Встановлено, що AI-технології можуть ефективно використовуватися на різних етапах навчального процесу: під час мотивації, пояснення нового матеріалу, закріплення знань та контролю навчальних досягнень.

Проведений аналіз показав, що генеративний AI та інтелектуальні навчальні системи сприяють персоналізації навчання, автоматизації рутинних педагогічних процесів та розвитку адаптивного освітнього середовища.

Разом із тим визначено необхідність формування критичного мислення учнів, педагогічного контролю використання AI та дотримання принципів академічної доброчесності.
```


% =========================
% ВИСНОВКИ
% =========================

\newpage

\section*{ВИСНОВКИ}
\addcontentsline{toc}{section}{ВИСНОВКИ}

У результаті дослідження встановлено, що технології штучного інтелекту мають значний потенціал для вдосконалення математичної освіти.

AI дозволяє реалізувати персоналізоване навчання, автоматизувати частину рутинної роботи вчителя та підвищити ефективність освітнього процесу.

% =========================
% ЗАГАЛЬНІ ВИСНОВКИ
% =========================

\section*{ЗАГАЛЬНІ ВИСНОВКИ}

\addcontentsline{toc}{section}{ЗАГАЛЬНІ ВИСНОВКИ}

У результаті виконання курсової роботи було досліджено особливості використання штучного інтелекту в професійній діяльності вчителя математики та визначено основні напрями інтеграції AI-технологій у математичну освіту.

Проведений аналіз наукових джерел засвідчив, що сучасні AI-технології активно трансформують освітній простір та створюють нові можливості для персоналізації навчання, автоматизації педагогічних процесів і підвищення ефективності освітньої діяльності. Встановлено, що найбільш перспективними напрямами використання AI в математичній освіті є генеративний штучний інтелект, інтелектуальні навчальні системи, адаптивні освітні платформи та автоматизоване оцінювання.

У ході дослідження було визначено, що генеративний AI може ефективно використовуватися для:
\begin{itemize}
\item створення математичних задач;
\item підготовки навчальних матеріалів;
\item автоматизованого формування тестів;
\item організації диференційованого навчання;
\item забезпечення швидкого зворотного зв’язку.
\end{itemize}

З’ясовано, що використання AI сприяє:
\begin{itemize}
\item персоналізації освітнього процесу;
\item підвищенню навчальної мотивації учнів;
\item економії часу педагога;
\item розвитку адаптивного навчального середовища.
\end{itemize}

Разом із тим встановлено, що впровадження AI супроводжується певними ризиками та обмеженнями. До основних проблем належать:
\begin{itemize}
\item можливість порушення академічної доброчесності;
\item залежність учнів від AI;
\item ризики використання недостовірної інформації;
\item недостатній рівень цифрової компетентності педагогів;
\item проблема цифрової нерівності.
\end{itemize}

У роботі доведено, що ефективність використання AI значною мірою залежить від педагогічної майстерності вчителя математики. Саме педагог повинен забезпечувати:
\begin{itemize}
\item контроль використання AI;
\item перевірку правильності AI-відповідей;
\item розвиток критичного мислення учнів;
\item дотримання принципів академічної доброчесності.
\end{itemize}

Під час дослідження було сформульовано методичні рекомендації щодо використання AI у професійній діяльності вчителя математики. Зокрема, визначено необхідність:
\begin{itemize}
\item педагогічно обґрунтованого використання AI;
\item розвитку цифрової компетентності педагогів;
\item формування навичок критичного аналізу AI-відповідей;
\item поєднання традиційних та цифрових методів навчання.
\end{itemize}

Отже, штучний інтелект має значний потенціал для вдосконалення математичної освіти та професійної діяльності вчителя математики. Водночас його використання повинно бути відповідальним, педагогічно доцільним та орієнтованим на розвиток особистості учня.

% =========================
% ВИКОРИСТАННЯ GENERATIVE AI
% =========================

\section*{Використання генеративного штучного інтелекту}

\addcontentsline{toc}{section}{Використання генеративного штучного інтелекту}

Під час підготовки курсової роботи використовувалися генеративні системи штучного інтелекту ChatGPT та Claude для формування структури дослідження, систематизації наукових джерел, редагування тексту та підготовки окремих фрагментів роботи. Усі наукові положення, висновки та теоретичні узагальнення базуються виключно на реальних наукових джерелах, проіндексованих у базі даних Scopus. Остаточне редагування, перевірка достовірності інформації та оформлення роботи виконані автором самостійно.

% =========================
% СПИСОК ВИКОРИСТАНИХ ДЖЕРЕЛ
% =========================

\newpage

\printbibliography

\end{document} Ось все що написано